\begin{table}[t]\centering\footnotesize
\caption{Draft label-based era classes for our three target houses. Printed codes are the OCR
targets for weeks 6--7 (SS = spring/summer, AW = autumn/winter).}
\label{tab:eras}
\setlength{\tabcolsep}{3pt}
\begin{tabular}{@{}>{\raggedright\arraybackslash}p{1.25cm}>{\raggedright\arraybackslash}p{3.15cm}>{\raggedright\arraybackslash}p{3.15cm}@{}}\toprule
House & Era classes (label cue) & Printed date code \\\midrule
CDG & pre-1988 (no AD year); 1988--2000 (AD year, hyphenated style code); 2001--07; 2008--17;
2018-- (season-letter code cycles) &
``AD''~+ production year on the care label (SS goods usually carry the previous year); style-code season
letter A\,=\,SS2008\dots T\,=\,AW2017, restarting at A\,=\,SS2018 \\\addlinespace
Margiela & 1988--96 (blank label, four white tacks); 1997--2008 (numbered 0--23 label; women's Line~1 stays blank until c.\,2008);
2009--14 (design team, ``Maison Martin Margiela''); 2015-- (``Maison Margiela'') &
season label ``Printemps-\'Et\'e / Automne-Hiver'' + year (to c.\,1993); care-tag COMM.\ number
encoding year and season (1999--2010); ``YYYY 01/02'' from AW2010 \\\addlinespace
Undercover & 1990--2002 (Tokyo era); 2003--14 (Paris era, men's line sold as UNDERCOVERISM);
2015--20 (unified label); 2021-- &
product code on care/price tags in seven successive formats, e.g.\ ``UC''~+ season letter
P\,=\,AW2015\dots Z\,=\,AW2020; from SS2021 ``UC''~+ 1/2 (SS/AW) + year letter (A\,=\,2021) \\
\bottomrule\end{tabular}\end{table}
